\chapter{Turingmaschinen}
\section{Turing-Maschinen}
\begin{figure}[H]
	\centering
	\begin{tikzpicture}
		\tikzstyle{every path}=[very thick]

		\edef\sizetape{0.7cm}
		\tikzstyle{tmtape}=[draw,minimum size=\sizetape]
		\tikzstyle{tmhead}=[arrow box,draw,minimum size=.5cm,arrow box
		arrows={east:.25cm, west:0.25cm}]

		%% Draw TM tape
		\begin{scope}[start chain=1 going right,node distance=-0.15mm]
			\node [on chain=1,tmtape,draw=none] {$\ldots $};
			\node [on chain=1,tmtape] {};
			\node [on chain=1,tmtape] (input) {b};
			\node [on chain=1,tmtape] {b};
			\node [on chain=1,tmtape] {a};
			\node [on chain=1,tmtape] {a};
			\node [on chain=1,tmtape] {a};
			\node [on chain=1,tmtape] {a};
			\node [on chain=1,tmtape] {};
			\node [on chain=1,tmtape,draw=none] {$\ldots $};
			\node [on chain=1] {\textbf{Input/Output Tape}};
		\end{scope}

		%% Draw TM Finite Control
		\begin{scope}
			[shift={(3cm,-5cm)},start chain=circle placed {at=(-\tikzchaincount*60:1.5)}]
			\foreach \i in {q_0,q_1,q_2,q_3,\ddots,q_n}
			\node [on chain] {$\i$};

			% Arrow to current state
			\node (center) {};
			\draw[->] (center) -- (circle-2);

			\node[rounded corners,draw=black,thick,fit=(circle-1) (circle-2) (circle-3)
				(circle-4) (circle-5) (circle-6),
				label=below:\textbf{Finite Control}] (fsbox)
			{};
		\end{scope}

		%% Draw TM head below (input) tape cell
		\node [tmhead,yshift=-.3cm] at (input.south) (head) {$q_1$};

		%% Link Finite Control with Head
		\path[->,draw] (fsbox.north) .. controls (4.5,-1) and (0,-2) .. node[right]
		(headlinetext)
		{}
		(head.south);
		\node[xshift=3cm] at (headlinetext)
		{\begin{tabular}{c}
				\textbf{Reading and Writing Head} \\
				\textbf{(moves in both directions)}
			\end{tabular}};

	\end{tikzpicture}
	\caption{Turing-Maschine}
	\label{fig:xdfdadf}
\end{figure}

\section{Halteproblem}
\begin{mydefinition}[Halteproblem]\label{mydefinition:halteproblem}
Das Problem, für jedes beliebge Programm $M$ und eine dazugehörige Eingabe $x$ zu entscheiden, ob $M$ auf $x$ hält oder nicht, wird \emph{Halteproblem} genannt.
\end{mydefinition}

\begin{mytheorem}[Halteproblem]
Das Halteproblem kann von keinem Algorithmus entschieden werden.
\end{mytheorem}

\begin{myremark}\label{myremark:H_lambda}
	Es lässt sich leicht zeigen, dass auch das Problem, für jedes beliebige Programm $M$ zu entscheiden, ob $M$ auf der Eingabe $\lambda$ (leeres Wort) hält, nicht berechenbar ist.
\end{myremark}


\section{Berechenbare Zahlen (Computable Numbers)}
\begin{mydefinition}[Berechenbare Zahl (informell)]\label{mydefinition:comp_numb}
	Eine reelle Zahl $x$ heisst berechenbar, wenn ein Algorithmus $A(n)$ existiert, der zu beliebigem natürlichen $n$ die Zahl $x$ auf mindestens $n$ Nachkommastellen genau berechnet.
\end{mydefinition}

\begin{myexample}[$\sqrt{2}$ ist berechenbar]
Im Folgenden skizzieren wir einen Algorithmus, der die Berechenbarkeit von $x:=\sqrt{2}$​ belegt. Am Beispiel der Genauigkeitsanforderung $n:=3$ demonstrieren wir das Verfahren der Intervallhalbierung, um die ersten drei Dezimalstellen von $\sqrt{2}$ (abgeschnitten) zu bestimmen.

Startintervall:
\[
\sqrt{2} \in \lrs{1,2},
\qquad
1^2 = 1 < 2,\quad 2^2 = 4 > 2.
\]

\textbf{1. Schritt}
\[
m=\frac{1+2}{2}=1.5,\qquad 1.5^2=2.25>2 \Rightarrow \sqrt{2}\in\lrs{1,1.5}.
\]

\textbf{2. Schritt}
\[
m=\frac{1+1.5}{2}=1.25,\qquad 1.25^2=1.5625<2 \Rightarrow \sqrt{2}\in\lrs{1.25,1.5}.
\]

\textbf{3. Schritt}
\[
m=\frac{1.25+1.5}{2}=1.375,\qquad 1.375^2=1.890625<2 \Rightarrow \sqrt{2}\in\lrs{1.375,1.5}.
\]

\textbf{4. Schritt}
\[
m=\frac{1.375+1.5}{2}=1.4375,\qquad 1.4375^2>2 \Rightarrow \sqrt{2}\in\lrs{1.375,1.4375}.
\]

Nach weiteren Halbierungen erhält man beispielsweise
\[
\sqrt{2}\in\lrs{1.4140625,\,1.41455078125}.
\]

Nun betrachten wir die abgeschnittenen Werte auf drei Dezimalstellen:
\[
\floor{1000\cdot 1.4140625} = 1414,
\qquad
\floor{1000\cdot 1.41455078125} = 1414.
\]

Da beide Intervallgrenzen denselben abgeschnittenen Wert liefern, folgt
\[
\floor{1000\sqrt{2}} = 1414.
\]

Damit ist
\[
\sqrt{2}=1.414\ldots
\]

und die auf drei Dezimalstellen abgeschnittene Darstellung lautet
\[
1.414.
\]
Ein entsprechendes Python-Programm ist in \cref{listing:n_digit_sqrt_2} in \cref{ch:appendix} zu finden.
\end{myexample}

\begin{myexample}
Für jede natürliche Zahl $n$ definieren wir die Funktion $f$ wie folgt: $f(n) = 1$, falls die $n$-te Turingmaschine auf dem Input $\lambda$ (leeres Wort) hält und $f(0) = 0$ sonst. Die reelle Zahl
\begin{align*}
	r := \sum_{n\in\N}f(n)\cdot 2^{-n}
\end{align*}
wird \emph{Haltezahl} genannt. Es ist klar, dass die Haltezahl nicht berechenbar ist. Könnten wir sie berechnen, könnten wir auch das Halteproblem entscheiden (siehe auch \cref{myremark:H_lambda}).
\end{myexample}